\part{Chimie (7 points)~: Acide azélaïque}
L'acide azélaïque est un composé organique de formule $\ce{C9H16O4}$, on le
trouve dans le blé, le seigle et l'orge. Il est un précurseur de divers produits
industriels et également un composant de certains produits de soin de la peau.

On utilisera la notation simplifiée $\ce{AH}$ pour cet acide et $\ce{A^{-}}$pour
sa base conjuguée.

Cet exercice vise~:
\begin{itemize}
    \item l'étude d'une solution aqueuse d'acide azélaïque~;
    \item la vérification du pourcentage d'acide azélaïque dans un produit
          cosmétique.
\end{itemize}

\section*{Partie 1~: Étude d'une solution aqueuse d'acide azélaïque}
On prépare une solution aqueuse d'acide azélaïque en dissolvant la masse
$m=\qty{0.188}{\g}$ de cet acide dans le volume $V=\qty{100}{\mL}$ d'eau
distillée. La mesure de $\mathrm{pH}$ de la solution à
$\qty{25}{\degreeCelsius}$ donne $\mathrm{pH}=3,28$.

\begin{donnees}
    \begin{itemize}
        \item $M(\ce{AH})=\qty{188}{\g\per\mol}$
    \end{itemize}
\end{donnees}

\begin{enumerate}
    \item Écrire l'équation qui modélise la réaction entre l'acide $\ce{AH}$ et
          l'eau.
    \item Dresser le tableau d'avancement de la réaction entre l'acide azélaïque
          et l'eau.
    \item Montrer que le taux d'avancement final de la réaction s'écrit
          $\tau=\frac{M\cdot V\cdot 10^{-\mathrm{pH}}}{m}$.
    \item Calculer la valeur de $\tau$. Déduire.
    \item Montrer que l'expression du quotient de la réaction $Q_{r,\eq}$ à
          l'équilibre s'écrit
          $Q_{r,\eq}=\frac{m\cdot\tau^{2}}{M\cdot V(1-\tau)}$.
    \item Vérifier que la valeur du $\mathrm{pK_{A}}$ du couple
          $(\ce{AH_{(aq)}}\ttslash{}\ce{A^-_{(aq)}})$ est
          $\mathrm{pK_{A}}=4,54$.
    \item Déterminer, parmi les espèces $\ce{AH_{(aq)}}$ et $\ce{A^-_{(aq)}}$,
          l'espèce qui prédomine dans la solution à l'équilibre.
\end{enumerate}



\section*{Partie 2~: Vérification du pourcentage d'acide azélaïque dans un produit cosmétique}
\begin{minipage}{0.78\linewidth}
    On dispose d'un produit cosmétique conçu pour le traitement de l'acné et de
    la rosacée en application topique. L'étiquette du produit porte l'indication
    \textbf{<<Acide Azélaïque 10\%>>} (Fig.~\ref{fig:1}). L'indication $10\%$
    désigne le pourcentage massique de l'acide Azélaïque dans le produit.

    À fin de vérifier l'indication précédente, on prépare une solution aqueuse
    $\mathrm{(S)}$ de concentration molaire $C$ et de volume $V=\qty{100}{\mL}$
    en dissolvant la masse ${m_{\text{(produit)}}=\qty{1.88}{\g}}$ du produit
    cosmétique dans l'eau distillée.

    On réalise le dosage du volume $V_{A}=\qty{10}{\mL}$ de la solution
    $\mathrm{(S)}$ par une solution aqueuse $\mathrm{(S_{B})}$ d'hydroxyde de
    sodium $\ce{Na^+_{(aq)} + HO^-_{(aq)}}$ de concentration molaire
    ${C_{B}=\qty{1e-2}{\mol\per\L}}$ en utilisant le montage de la
    Fig.~\ref{fig:2}.  La courbe de la Fig.~\ref{fig:3} représente les
    variations du $\mathrm{pH}$ du mélange en fonction du volume $V_{B}$ de la
    solution $\mathrm{(S_{B})}$ ajoutée.
\end{minipage}
\hfill
\begin{minipage}{0.20\linewidth}
    \begin{figure}[H]
        \centering
        \includegraphics[width=\textwidth]{2bac_svt_national_2024_rattrapage_chim_1}
        \caption{}
        \label{fig:2bac_svt_national_2024_rattrapage_chim_1}
    \end{figure}
\end{minipage}

\begin{minipage}{0.42\linewidth}
    \begin{figure}[H]
        \centering
        \includegraphics[width=\textwidth]{2bac_svt_national_2024_rattrapage_chim_2}
        \caption{}
        \label{fig:2bac_svt_national_2024_rattrapage_chim_2}
    \end{figure}
\end{minipage}
\hfill
\begin{minipage}{0.42\linewidth}
    \begin{figure}[H]
        \centering
        \includegraphics[width=\textwidth]{2bac_svt_national_2024_rattrapage_chim_3}
        \caption{}
        \label{fig:2bac_svt_national_2024_rattrapage_chim_3}
    \end{figure}
\end{minipage}

\begin{enumerate}
    \item Attribuer au numéros~(1),~(2) et~(3) du montage de la Fig.~\ref{fig:2}
          les noms correspondants.
    \item Écrire l'équation de la réaction du dosage supposée totale.
    \item Déterminer les coordonnées $(V_{B\mathrm{E}}~;~\mathrm{pH_{E}})$ du
          point d'équivalence.
    \item Calculer la concentration $C$ de la solution $\mathrm{(S)}$.
    \item Vérifier l'indication <<Acide Azélaïque $10\%$>> indiquée sur
          l'étiquette du produit cosmétique.
\end{enumerate}




